\part*{Guide: doing the TDK}\addcontentsline{toc}{part}{Guide: doing the TDK}
%----------------------------------------------------------------------------
\chapter{Start here}\label{app:here}
%----------------------------------------------------------------------------
This guide explains how to start, write and submit a TDK paper with FTSRG advisors Martin Farkas (\href{https://github.com/BlackLight54}{\texttt{BlackLight54}} on GitHub) and Imre Kocsis (\href{https://github.com/imrekocsis}{\texttt{imrekocsis}}). It covers our working process and recommended tools.
The guide follows the appendices in the paper template. You do not need to delete it before submission. Build \texttt{reviewer.tex} and upload the resulting \texttt{reviewer.pdf}, which excludes the guide and all notes (\cref{app:latex:todo}).

The template uses three kinds of notes. \todo[inline]{An orange note is for you: what a part of the paper must contain.}
\advisortodo{A blue note is for us: a draft we still have to confirm, or a link we still have to add.}
\atodo[inline]{A green note, marked A, is our review comment on your paper. Resolve it, then delete it.}

\paragraph{How the guide states rules.}\label{app:here:words} In this guide, \emph{must} marks a requirement and always names its source: the faculty (the call, the review sheet or the faculty's AI guideline) or us, the advisors. \emph{We recommend} marks our default, which you may depart from with a reason. \emph{You can} marks an option.

\section{Where to look}\label{app:here:route}
\begin{description}
  \item[Your first week] Work through \cref{app:start} in order.
  \item[This week] Find your phase in \cref{tab:lifecycle}, and bring your progress to the weekly consultation (\cref{app:consult}).
  \item[Writing a chapter] \Cref{app:chapters} says what each chapter contains, with a worked example. \Cref{app:writing:essentials} lists the elements the lint measures, and \cref{app:latex} shows the syntax.
  \item[Before upload] Set \texttt{\textbackslash lintphase} to \texttt{submit}, run \texttt{make check-submit} and go through \cref{tab:polish}. The deadlines are in \cref{app:tdk:dates}.
  \item[After upload] \Cref{tab:lifecycle} covers the review weeks and the conference, and \cref{app:lifecycle:presentation} the presentation.
  \item[A word you do not know] \Cref{app:glossary} defines the process terms this guide leans on: anchor, spine, beats, essentials, hill climb, premortem, skim view.
  \item[Working with an agent] \texttt{AGENTS.md}, at the paper repository's root, tells an agent how to build, check and change the paper, and points into this guide.
\end{description}
